\section{Introduction}
An animal searching for a goal can use recent spatial experience to avoid repeatedly exploring the same locations. A short-term record of recent visits can bias the next movement while still allowing a return when the route requires it. This record must persist beyond the neural activity associated with an individual observation and remain accessible to the circuits selecting subsequent actions. Short-term synaptic facilitation provides one neuronal mechanism for retaining information beyond the initiating spikes \cite{mongillo2008synaptic}. 

Astrocytes are a candidate substrate because they respond to neuronal activity and can regulate the neurons with which they interact. Experimental activation of hippocampal astrocytes can induce synaptic potentiation and improve memory performance \cite{adamsky2018astrocytic}. Astrocytic signaling also regulates dendritic integration and contributes to spatial learning \cite{bohmbach2022astrocytic}. These findings connect astrocyte function to the processing and retention of behaviorally relevant information. They motivate investigating how an astrocytic state could carry information from a recent observation into a later behavioral choice.

Calcium-dependent modulation offers a mechanism for this persistence. A calcium response that outlasts its initiating input can continue to affect neuronal activity after that input has ended. Computational models have used astrocyte-dependent modulation to maintain and retrieve neuronal patterns, supporting a possible role in working memory \cite{gordleeva2021working,depitta2022multiple}. Inhibition provides a route by which such a memory could influence action selection. Extracellular ATP can excite hippocampal interneurons, and neuron--glia signaling can increase inhibitory transmission \cite{bowser2004atp,torres2012extracellular}. We hypothesize that coupling a persistent astrocytic state to inhibition can bias choices according to recent spatial experience.

A direct computational precedent is the navigation model of Tsybina and colleagues \cite{tsybina2026}. Their spiking neuron--astrocyte network uses calcium-dependent inhibition to suppress recently visited locations, alongside synaptic plasticity that reinforces successful behavior. The present work examines how this memory mechanism can interact with learned predictions of action outcomes. Before a movement is executed, its expected outcome can be compared with a record of recent visits. This makes it possible to use memory prospectively, inhibiting actions expected to lead back to familiar regions.

We introduce AstroRL, which combines a recurrent spiking network, an action-conditioned transition predictor, an astrocyte-inspired calcium memory, and a spiking motor circuit. Actual observations activate overlapping spatial features and update a reduced glutamate--calcium process. Predictions of candidate outcomes query the resulting memory, and greater predicted familiarity recruits stronger motor inhibition. The feedback changes action probabilities while preserving the possibility of revisits. Prediction error trains the recurrent representation through e-prop \cite{bellec2020}, while reward trains a separate, detached motor readout. This architecture allows the prediction system to continue learning as the policy changes, without a direct reward gradient into the recurrent network.

We evaluate exploration and goal-directed navigation on fixed, position-observed maps. Exploration with a fixed motor readout tests whether familiarity memory can change behavior before reward-based policy learning. Reinforcement-learning experiments then assess its contribution to a trained motor policy. Removing memory at evaluation, with learned weights fixed, tests whether that policy continues to use the memory signal. An exponential neuronal trace provides a functional comparison, while shared-input probes evaluate transition prediction independently of the trajectories selected by each policy.

Our contribution is a prediction-guided interface between astrocyte-inspired memory and action selection, evaluated with and without motor learning. On the tested maps, calcium memory increases coverage without policy learning and supports goal-directed navigation in structured environments. Its performance remains competitive with the neuronal-trace control, and removing memory reduces the performance of trained policies. These findings provide computational support for the hypothesis that astrocytic dynamics could fulfill trace-like memory functions in biological circuits.